\section{Timeline Grammar: Domain IR}
\label{sec:timeline-grammar}

The Timeline Grammar is the first grammar implemented in the diagram compiler. This section specifies its Domain IR---the declarative representation of \emph{what} a timeline communicates.

\textbf{Note:} This section re-homes and re-contextualizes the original Timeline IR specification as ``Grammar \#1'' in the broader diagram compiler architecture. The IR definition is unchanged; only the framing is updated.

\subsection{Design Principles}

The Timeline Domain IR follows the general grammar design principles (Section~\ref{sec:grammar-concept}), with timeline-specific elaboration:

\begin{enumerate}
  \item \textbf{Make illegal states unrepresentable.} The IR's type system should prevent
        malformed timelines at the structural level, not merely validate them after the fact.
  \item \textbf{Rendering-agnostic.} The IR describes meaning, not pixels. Visual
        decisions---colors, fonts, positioning---belong to the theme engine.
  \item \textbf{Git-friendly.} Stable field ordering, line-oriented YAML syntax, deterministic
        serialization, minimal diff churn.
  \item \textbf{Agent-generatable.} Clear required-vs-optional distinctions, explicit defaults,
        forgiving structure that remains validatable.
  \item \textbf{Orthogonal core.} A small set of primitives; convenience features are optional
        layers, not core complexity.
\end{enumerate}

\subsection{Type Vocabulary}
\label{sec:ir-types}

The IR uses the following type vocabulary consistently throughout this specification:

\begin{table}[h]
\centering
\begin{tabular}{lp{10cm}}
\toprule
\textbf{Type} & \textbf{Description} \\
\midrule
\texttt{string} & UTF-8 text, non-empty unless explicitly noted \\
\texttt{string?} & Optional string (may be omitted or null) \\
\texttt{int} & Integer \\
\texttt{number} & Decimal number \\
\texttt{number[0..1]} & Decimal in closed interval $[0, 1]$ \\
\texttt{bool} & Boolean (\texttt{true}/\texttt{false}) \\
\texttt{date} & A temporal point (see \S\ref{sec:date-model}) \\
\texttt{daterange} & A temporal interval with start and optional end \\
\texttt{id} & Document-unique identifier (kebab-case slug, e.g., \texttt{alpha-release}) \\
\texttt{ref<T>} & Reference to an entity of type \texttt{T} by its \texttt{id} \\
\texttt{list<T>} & Ordered list of elements of type \texttt{T} \\
\texttt{optional<T>} & Value of type \texttt{T} that may be omitted \\
\texttt{enum\{...\}} & One of the listed literal values \\
\texttt{map<K,V>} & Key-value mapping \\
\bottomrule
\end{tabular}
\caption{IR type vocabulary}
\label{tab:type-vocab}
\end{table}

\subsection{Document Structure}
\label{sec:ir-structure}

A Timeline IR document is a single YAML file with the following root structure:

\begin{lstlisting}[language=yaml,caption={Root document structure}]
version: "1.0"
metadata: { ... }
tracks: [ ... ]
groups: [ ... ]       # optional
activities: [ ... ]
milestones: [ ... ]   # optional
annotations: [ ... ]  # optional
sections: [ ... ]     # optional
legend: { ... }       # optional
\end{lstlisting}

\subsubsection{Root Fields}

\begin{table}[h]
\centering
\small
\begin{tabular}{llccp{5cm}}
\toprule
\textbf{Field} & \textbf{Type} & \textbf{Req} & \textbf{Default} & \textbf{Description} \\
\midrule
\texttt{version} & \texttt{string} & Yes & --- & Specification version (semver, e.g., \texttt{"1.0"}) \\
\texttt{metadata} & \texttt{Metadata} & Yes & --- & Document-level metadata \\
\texttt{tracks} & \texttt{list<Track>} & Yes & --- & Swimlanes/rows; at least one required \\
\texttt{groups} & \texttt{list<Group>} & No & \texttt{[]} & Logical groupings of activities \\
\texttt{activities} & \texttt{list<Activity>} & Yes & --- & Time-spanning items; may be empty \\
\texttt{milestones} & \texttt{list<Milestone>} & No & \texttt{[]} & Point-in-time markers \\
\texttt{annotations} & \texttt{list<Annotation>} & No & \texttt{[]} & Callouts, notes, highlights \\
\texttt{sections} & \texttt{list<Section>} & No & \texttt{[]} & Visual/logical document sections \\
\texttt{legend} & \texttt{Legend} & No & auto & Legend mapping; auto-generated if omitted \\
\bottomrule
\end{tabular}
\caption{Root document fields}
\label{tab:root-fields}
\end{table}

\textbf{Invariant:} The document must contain at least one track. Activities and milestones
may both be empty, but a timeline with neither is semantically vacuous (valid but degenerate).


%% ============================================================================
%% METADATA
%% ============================================================================
\subsection{Metadata}
\label{sec:ir-metadata}

The \texttt{metadata} object contains document-level information and temporal framing.

\begin{table}[h]
\centering
\small
\begin{tabular}{llccp{4.5cm}}
\toprule
\textbf{Field} & \textbf{Type} & \textbf{Req} & \textbf{Default} & \textbf{Description} \\
\midrule
\texttt{title} & \texttt{string} & Yes & --- & Primary document title \\
\texttt{subtitle} & \texttt{string?} & No & \texttt{null} & Secondary title or tagline \\
\texttt{author} & \texttt{string?} & No & \texttt{null} & Author or team name \\
\texttt{created} & \texttt{date} & No & now & Document creation date \\
\texttt{updated} & \texttt{date} & No & now & Last modification date \\
\texttt{time\_range} & \texttt{TimeRange} & Yes & --- & Temporal bounds of the timeline \\
\texttt{axis\_unit} & \texttt{AxisUnit} & No & inferred & Granularity of the time axis \\
\texttt{theme} & \texttt{string?} & No & \texttt{null} & Theme identifier (resolved by renderer) \\
\texttt{locale} & \texttt{string?} & No & \texttt{"en"} & BCP-47 locale for date formatting \\
\texttt{today} & \texttt{date?} & No & eval time & Reference date for \texttt{now}, relative dates, and today-marker; see \S\ref{sec:date-anchor} \\
\texttt{fiscal\_year\_start} & \texttt{int [1..12]} & No & \texttt{1} & Calendar month on which the fiscal year begins (1~=~January); see \S\ref{sec:date-anchor} \\
\texttt{description} & \texttt{string?} & No & \texttt{null} & Extended description or notes \\
\bottomrule
\end{tabular}
\caption{Metadata fields}
\label{tab:metadata-fields}
\end{table}

\subsubsection{TimeRange}

The \texttt{time\_range} defines the temporal viewport of the timeline---the span that will
be rendered. Activities or milestones outside this range are valid but may be clipped.

\begin{lstlisting}[language=yaml]
time_range:
  start: 2026-Q1
  end: 2026-Q4
\end{lstlisting}

\begin{table}[h]
\centering
\begin{tabular}{llccp{5cm}}
\toprule
\textbf{Field} & \textbf{Type} & \textbf{Req} & \textbf{Default} & \textbf{Description} \\
\midrule
\texttt{start} & \texttt{date} & Yes & --- & Inclusive start of the viewport \\
\texttt{end} & \texttt{date} & No & open & Inclusive end; omit for open-ended \\
\bottomrule
\end{tabular}
\caption{TimeRange fields}
\label{tab:timerange-fields}
\end{table}

\subsubsection{AxisUnit}

\begin{lstlisting}
enum AxisUnit { day, week, month, quarter, half, year }
\end{lstlisting}

If omitted, the renderer infers axis unit from the time range span. The axis unit affects
visual tick marks and labels, not the precision of individual dates.

\subsubsection{Date Anchor and Fiscal Calendar}
\label{sec:date-anchor}

\paragraph{Date anchor (\texttt{today}).}
The symbolic date \texttt{now} and all relative date offsets (\texttt{+3m}, \texttt{-2w},
etc.) resolve against the \emph{date anchor}, evaluated in this order:
\begin{enumerate}
  \item \texttt{metadata.today} — if present, this value is used exclusively.
  \item \texttt{metadata.created} — fallback when \texttt{today} is absent.
  \item \textbf{Hard error} — if neither is present and any \texttt{now} or relative date
        appears in the document, the document is not deterministically evaluable and
        validation must reject it.
\end{enumerate}

Renderers \textbf{must not} consult the system clock to resolve \texttt{now} or relative
dates. The today-marker annotation (\texttt{type: today-marker}) also uses this anchor.

\textbf{Determinism caveat:} When \texttt{today} is omitted and \texttt{created} serves
as the anchor, output is deterministic (the \texttt{created} date is fixed in the
document). However, the document's visual "current position" does not advance with real
time. Authors and agents \textbf{should} set \texttt{today} explicitly whenever \texttt{now}
or relative dates appear, so that the document's meaning is unambiguous and byte-stable
across environments and rendering runs.

\paragraph{Fiscal calendar (\texttt{fiscal\_year\_start}).}
\texttt{fiscal\_year\_start} specifies the calendar month (1--12) on which the fiscal year
begins. The default is~1 (January), which makes fiscal and calendar years identical.

Fiscal quarter dates (\texttt{FYnn-Qk}) map to calendar dates as follows. \texttt{FY}$nn$
denotes the fiscal year beginning on month \texttt{fiscal\_year\_start} of calendar year
$20nn$. Fiscal quarter $k$ ($k \in \{1,2,3,4\}$) spans three consecutive calendar months
starting at calendar month $m_k$, where:
\[
  m_k = ((\texttt{fiscal\_year\_start} - 1) + (k - 1) \times 3) \bmod 12 + 1
\]
and the calendar year is $20nn$ if $m_k \ge \texttt{fiscal\_year\_start}$, otherwise
$20nn + 1$.

\textbf{Examples with \texttt{fiscal\_year\_start: 1} (default):}
FY26-Q1 = Jan--Mar 2026; FY26-Q2 = Apr--Jun 2026 (identical to calendar quarters).

\textbf{Examples with \texttt{fiscal\_year\_start: 4} (April fiscal year):}
FY26-Q1 = Apr--Jun 2026; FY26-Q2 = Jul--Sep 2026; FY26-Q3 = Oct--Dec 2026;
FY26-Q4 = Jan--Mar 2027.

If any \texttt{FY...} date appears in the document and \texttt{fiscal\_year\_start}
$\neq 1$, authors \textbf{should} set \texttt{fiscal\_year\_start} explicitly to ensure
all renderers produce consistent x-coordinates for fiscal dates.


%% ============================================================================
%% DATE MODEL
%% ============================================================================
\subsection{Date Model}
\label{sec:date-model}

Timelines present temporal information at varying granularities. A product roadmap may show
quarters; a conference schedule shows hours. The IR must represent this heterogeneity without
forcing false precision or losing meaningful vagueness.

\subsubsection{Date Representations}

A \texttt{date} in the IR is one of the following forms:

\begin{table}[h]
\centering
\small
\begin{tabular}{llp{6cm}}
\toprule
\textbf{Form} & \textbf{Example} & \textbf{Semantics} \\
\midrule
ISO date & \texttt{2026-06-09} & Specific calendar day \\
ISO datetime & \texttt{2026-06-09T14:00} & Specific moment (minute precision) \\
Year-month & \texttt{2026-06} & Entire month (June 2026) \\
Year only & \texttt{2026} & Entire year \\
Quarter & \texttt{2026-Q2} & Calendar quarter (Apr--Jun) \\
Half & \texttt{2026-H1} & Calendar half (Jan--Jun) \\
Fiscal quarter & \texttt{FY26-Q2} & Fiscal quarter (calendar mapping per \texttt{fiscal\_year\_start}; see \S\ref{sec:date-anchor}) \\
Relative & \texttt{+3m} & 3 months from the date anchor (\texttt{metadata.today} $\to$ \texttt{metadata.created}; see \S\ref{sec:date-anchor}) \\
Symbolic & \texttt{now} & Resolves to the date anchor (\texttt{metadata.today} $\to$ \texttt{metadata.created}; see \S\ref{sec:date-anchor}) \\
\bottomrule
\end{tabular}
\caption{Date representation forms}
\label{tab:date-forms}
\end{table}

\textbf{Design rationale:} We deliberately support coarse granularity (quarters, halves, years)
as first-class concepts rather than forcing conversion to day-level dates. This preserves
authorial intent: ``Q3 2026'' means the entire quarter, not ``2026-07-01''.

\subsubsection{Uncertain and Open Dates}

For dates that are genuinely unknown, tentative, or open-ended, the IR provides explicit
encodings rather than using null-as-magic:

\begin{table}[h]
\centering
\begin{tabular}{lp{8cm}}
\toprule
\textbf{Value} & \textbf{Semantics} \\
\midrule
\texttt{tbd} & Date is not yet determined; renders as ``TBD'' \\
\texttt{ongoing} & No planned end; activity continues indefinitely \\
\texttt{unknown} & Date exists but is not known to the author \\
\texttt{\textasciitilde 2026-Q3} & Approximate/tentative (prefix tilde) \\
\bottomrule
\end{tabular}
\caption{Uncertain and open date markers}
\label{tab:uncertain-dates}
\end{table}

\textbf{Invariant:} \texttt{tbd}, \texttt{ongoing}, and \texttt{unknown} are valid only in
positions where the schema allows uncertainty (e.g., \texttt{end} of a daterange, not
\texttt{start} of the document's \texttt{time\_range}).

\subsubsection{DateRange}

A \texttt{daterange} represents a temporal interval:

\begin{lstlisting}[language=yaml]
# Explicit start and end
start: 2026-04-01
end: 2026-06-30

# Open-ended (no planned end)
start: 2026-04-01
end: ongoing

# Shorthand for single-granularity span
span: 2026-Q2    # equivalent to start: 2026-04-01, end: 2026-06-30
\end{lstlisting}

\begin{table}[h]
\centering
\begin{tabular}{llccp{5cm}}
\toprule
\textbf{Field} & \textbf{Type} & \textbf{Req} & \textbf{Default} & \textbf{Description} \\
\midrule
\texttt{start} & \texttt{date} & Yes* & --- & Interval start (inclusive) \\
\texttt{end} & \texttt{date|ongoing|tbd} & No & \texttt{ongoing} & Interval end (inclusive); omitting is equivalent to \texttt{end: ongoing} \\
\texttt{span} & \texttt{date} & Yes* & --- & Single-unit shorthand; mutually exclusive with \texttt{start}/\texttt{end} \\
\bottomrule
\end{tabular}
\caption{DateRange fields (*exactly one of \texttt{span} or \texttt{start} (+ optional \texttt{end}) is required; co-presence of \texttt{span} with \texttt{start} or \texttt{end} is a well-formedness error)}
\label{tab:daterange-fields}
\end{table}

\textbf{Invariant:} When both \texttt{start} and \texttt{end} are concrete dates,
\texttt{end >= start}. The \texttt{span} shorthand expands to the natural bounds of its
granularity (e.g., \texttt{span: 2026-Q2} expands to Q2's first and last days).

\textbf{Invariant (exclusivity):} \texttt{span} is \emph{mutually exclusive} with
\texttt{start} and \texttt{end}. A daterange must use exactly one of:
\begin{itemize}
  \item \texttt{span} alone (shorthand form), or
  \item \texttt{start} alone or \texttt{start} + \texttt{end} (explicit form).
\end{itemize}
Co-presence of \texttt{span} with \texttt{start} or \texttt{end} on the same entity is
a well-formedness error (\texttt{SPAN\_START\_CONFLICT}).

\textbf{Open/ongoing semantics:} An Activity or daterange with \texttt{start} specified
but \texttt{end} omitted is an \emph{open/ongoing interval}. This is semantically
identical to writing \texttt{end: ongoing} and is an explicit, valid state (not an
authoring error). Renderers must extend the bar to the canvas right edge and apply
an open-end indicator (e.g., right-pointing chevron).


%% ============================================================================
%% TRACKS
%% ============================================================================
\subsection{Tracks (Swimlanes)}
\label{sec:ir-tracks}

Tracks are horizontal lanes that organize activities visually. Every activity belongs to
exactly one track. Tracks appear in document order (first track at top).

\begin{lstlisting}[language=yaml,caption={Track examples}]
tracks:
  - id: platform
    label: Platform Team
    description: Core infrastructure work

  - id: mobile
    label: Mobile Apps
    color: blue         # optional hint
\end{lstlisting}

\begin{table}[h]
\centering
\small
\begin{tabular}{llccp{4.5cm}}
\toprule
\textbf{Field} & \textbf{Type} & \textbf{Req} & \textbf{Default} & \textbf{Description} \\
\midrule
\texttt{id} & \texttt{id} & Yes & --- & Unique track identifier \\
\texttt{label} & \texttt{string} & Yes & --- & Display label \\
\texttt{description} & \texttt{string?} & No & \texttt{null} & Extended description \\
\texttt{color} & \texttt{string?} & No & theme & Color hint (theme may override) \\
\texttt{group} & \texttt{ref<Group>?} & No & \texttt{null} & Parent group (for nested tracks) \\
\texttt{index} & \texttt{int?} & No & position & Explicit sort index; unique non-negative integer; tracks render top-to-bottom in ascending \texttt{index} order \\
\texttt{collapsed} & \texttt{bool} & No & \texttt{false} & Initial collapsed state (if renderer supports) \\
\bottomrule
\end{tabular}
\caption{Track fields}
\label{tab:track-fields}
\end{table}

\textbf{Invariant:} Track \texttt{id} values must be unique within the document.

%% ============================================================================
%% GROUPS
%% ============================================================================
\subsection{Groups}
\label{sec:ir-groups}

Groups provide logical clustering of activities, tracks, or other groups. They enable
hierarchical organization (e.g., ``Engineering'' containing ``Platform'' and ``Mobile'' tracks)
without conflating logical grouping with visual swimlanes.

\begin{lstlisting}[language=yaml,caption={Group example}]
groups:
  - id: engineering
    label: Engineering
    children:
      - platform    # ref to track
      - mobile      # ref to track

  - id: releases
    label: Release Milestones
    members:
      - alpha-release   # ref to milestone
      - beta-release
\end{lstlisting}

\begin{table}[h]
\centering
\small
\begin{tabular}{llccp{4.5cm}}
\toprule
\textbf{Field} & \textbf{Type} & \textbf{Req} & \textbf{Default} & \textbf{Description} \\
\midrule
\texttt{id} & \texttt{id} & Yes & --- & Unique group identifier \\
\texttt{label} & \texttt{string} & Yes & --- & Display label \\
\texttt{description} & \texttt{string?} & No & \texttt{null} & Extended description \\
\texttt{parent} & \texttt{ref<Group>?} & No & \texttt{null} & Parent group (for nesting) \\
\texttt{children} & \texttt{list<ref<Track|Group>>} & No & \texttt{[]} & Ordered child tracks/groups \\
\texttt{members} & \texttt{list<ref<Activity|Milestone>>} & No & \texttt{[]} & Member entities \\
\texttt{color} & \texttt{string?} & No & theme & Color hint \\
\bottomrule
\end{tabular}
\caption{Group fields}
\label{tab:group-fields}
\end{table}

\textbf{Invariant:} Group hierarchies must be acyclic. A group cannot be its own ancestor.


%% ============================================================================
%% ACTIVITIES
%% ============================================================================
\subsection{Activities}
\label{sec:ir-activities}

Activities are time-spanning items---phases, projects, tasks, or efforts. They are the
primary content of most timelines.

\begin{lstlisting}[language=yaml,caption={Activity examples}]
activities:
  - id: api-redesign
    label: API Redesign
    track: platform
    start: 2026-Q1
    end: 2026-Q2
    status: in-progress
    progress: 0.4

  - id: mobile-launch
    label: Mobile App Launch
    track: mobile
    span: 2026-Q3
    status: planned
    category: launch
\end{lstlisting}

\begin{table}[h]
\centering
\small
\begin{tabular}{llccp{4cm}}
\toprule
\textbf{Field} & \textbf{Type} & \textbf{Req} & \textbf{Default} & \textbf{Description} \\
\midrule
\texttt{id} & \texttt{id} & Yes & --- & Unique activity identifier \\
\texttt{label} & \texttt{string} & Yes & --- & Display label (shown on bar) \\
\texttt{track} & \texttt{ref<Track>} & Yes & --- & Containing track \\
\texttt{start} & \texttt{date} & Yes* & --- & Start date \\
\texttt{end} & \texttt{date|ongoing|tbd} & No & \texttt{ongoing} & End date; omitting is equivalent to \texttt{end: ongoing} (open interval) \\
\texttt{span} & \texttt{date} & Yes* & --- & Single-unit shorthand; mutually exclusive with \texttt{start}/\texttt{end} \\
\texttt{status} & \texttt{Status} & No & \texttt{planned} & Visual communication status \\
\texttt{progress} & \texttt{number[0..1]?} & No & \texttt{null} & Completion fraction \\
\texttt{category} & \texttt{string?} & No & \texttt{null} & Semantic category (for styling) \\
\texttt{description} & \texttt{string?} & No & \texttt{null} & Extended description \\
\texttt{group} & \texttt{ref<Group>?} & No & \texttt{null} & Logical group membership \\
\texttt{url} & \texttt{string?} & No & \texttt{null} & External link \\
\texttt{tags} & \texttt{list<string>} & No & \texttt{[]} & Freeform tags \\
\texttt{metadata} & \texttt{map<string,any>} & No & \texttt{\{\}} & Extension metadata \\
\bottomrule
\end{tabular}
\caption{Activity fields (*exactly one of \texttt{span} or \texttt{start} (+ optional \texttt{end}) is required; co-presence of \texttt{span} with \texttt{start} or \texttt{end} is a well-formedness error)}
\label{tab:activity-fields}
\end{table}

\subsubsection{Status Enum}
\label{sec:status-enum}

The \texttt{status} field communicates \emph{visual} state for presentation purposes.
This is explicitly \textbf{not} a project-management workflow state; it describes what
the audience should understand about the item's current condition.

\begin{lstlisting}
enum Status {
  planned,      # Scheduled but not yet started
  in-progress,  # Currently active
  done,         # Completed
  at-risk,      # May miss target (visual warning)
  blocked,      # Cannot proceed (visual alert)
  cancelled,    # No longer planned
  tentative     # Not yet confirmed
}
\end{lstlisting}

\textbf{Design rationale:} We include \texttt{at-risk} and \texttt{blocked} because executive
timelines frequently need to communicate risk visually. We omit workflow states like
``in review'' or ``approved'' which belong to project-management tools, not visual
communication artifacts.

\subsubsection{Progress}

The \texttt{progress} field is a number in $[0, 1]$ representing completion fraction.
It has no inherent visual meaning---the theme decides whether to render it as a
progress bar, percentage label, or ignore it entirely.

\textbf{Invariant:} If \texttt{progress} is provided, it must be in $[0, 1]$.
\texttt{progress: 1.0} does not imply \texttt{status: done}; these are orthogonal.


%% ============================================================================
%% MILESTONES
%% ============================================================================
\subsection{Milestones}
\label{sec:ir-milestones}

Milestones are point-in-time markers---releases, deadlines, decision points, or events.
Unlike activities, they have no duration.

\begin{lstlisting}[language=yaml,caption={Milestone examples}]
milestones:
  - id: alpha-release
    label: Alpha Release
    date: 2026-03-15
    track: platform
    status: done
    icon: flag
    color: cyan

  - id: board-review
    label: Board Review
    date: 2026-Q2        # first day of Q2
    category: governance
\end{lstlisting}

\begin{table}[h]
\centering
\small
\begin{tabular}{llccp{4.5cm}}
\toprule
\textbf{Field} & \textbf{Type} & \textbf{Req} & \textbf{Default} & \textbf{Description} \\
\midrule
\texttt{id} & \texttt{id} & Yes & --- & Unique milestone identifier \\
\texttt{label} & \texttt{string} & Yes & --- & Display label \\
\texttt{date} & \texttt{date} & Yes & --- & Point in time \\
\texttt{track} & \texttt{ref<Track>?} & No & global & Associated track (or global) \\
\texttt{status} & \texttt{Status} & No & \texttt{planned} & Visual status \\
\texttt{category} & \texttt{string?} & No & \texttt{null} & Semantic category \\
\texttt{icon} & \texttt{string?} & No & theme & Icon hint (diamond, flag, star, etc.) \\
\texttt{color} & \texttt{string?} & No & theme & Color hint (theme may override) \\
\texttt{description} & \texttt{string?} & No & \texttt{null} & Extended description \\
\texttt{group} & \texttt{ref<Group>?} & No & \texttt{null} & Logical group membership \\
\texttt{url} & \texttt{string?} & No & \texttt{null} & External link \\
\texttt{tags} & \texttt{list<string>} & No & \texttt{[]} & Freeform tags \\
\texttt{metadata} & \texttt{map<string,any>} & No & \texttt{\{\}} & Application data; renderers ignore unknown keys \\
\bottomrule
\end{tabular}
\caption{Milestone fields}
\label{tab:milestone-fields}
\end{table}

\textbf{Note:} When \texttt{track} is omitted, the milestone is ``global'' and may be
rendered spanning all tracks (e.g., a vertical line with label at top).

%% ============================================================================
%% ANNOTATIONS
%% ============================================================================
\subsection{Annotations}
\label{sec:ir-annotations}

Annotations add contextual information---callouts, notes, highlights, or today-markers---
without being first-class timeline entities.

\begin{lstlisting}[language=yaml,caption={Annotation examples}]
annotations:
  - type: today-marker
    date: now

  - type: callout
    target: api-redesign
    text: "Critical path item"
    position: above

  - type: bracket
    targets: [alpha-release, beta-release]
    label: "Testing Phase"

  - type: period
    start: 2026-06-01
    end: 2026-06-15
    label: "Code Freeze"
    style: highlight
\end{lstlisting}

\begin{table}[h]
\centering
\small
\begin{tabular}{llccp{4cm}}
\toprule
\textbf{Field} & \textbf{Type} & \textbf{Req} & \textbf{Default} & \textbf{Description} \\
\midrule
\texttt{type} & \texttt{AnnotationType} & Yes & --- & Annotation kind \\
\texttt{id} & \texttt{id?} & No & auto & Optional identifier \\
\texttt{target} & \texttt{ref<Activity|Milestone>?} & Cond & --- & Single target entity \\
\texttt{targets} & \texttt{list<ref<...>>?} & Cond & --- & Multiple targets \\
\texttt{date} & \texttt{date?} & Cond & --- & Point annotation \\
\texttt{start}/\texttt{end} & \texttt{date?} & Cond & --- & Range annotation \\
\texttt{text}/\texttt{label} & \texttt{string?} & No & \texttt{null} & Display text \\
\texttt{position} & \texttt{enum\{above,below,left,right\}?} & No & auto & Hint only \\
\texttt{style} & \texttt{string?} & No & theme & Style hint \\
\bottomrule
\end{tabular}
\caption{Annotation fields (conditional fields depend on type)}
\label{tab:annotation-fields}
\end{table}

\subsubsection{Annotation Types}

\begin{lstlisting}
enum AnnotationType {
  today-marker,  # Vertical line at current date
  callout,       # Note attached to an entity
  bracket,       # Visual bracket spanning entities
  period,        # Highlighted time period
  note,          # Freestanding note
  connector      # Line connecting two entities
}
\end{lstlisting}


%% ============================================================================
%% SECTIONS
%% ============================================================================
\subsection{Sections}
\label{sec:ir-sections}

Sections divide the timeline into logical or visual chapters. They enable multi-page
timelines or distinct phases that should be visually separated.

\begin{lstlisting}[language=yaml,caption={Section example}]
sections:
  - id: current-state
    label: "Current State (Q1-Q2)"
    time_range:
      start: 2026-Q1
      end: 2026-Q2

  - id: future-state
    label: "Future Roadmap (Q3-Q4)"
    time_range:
      start: 2026-Q3
      end: 2026-Q4
\end{lstlisting}

\begin{table}[h]
\centering
\small
\begin{tabular}{llccp{5cm}}
\toprule
\textbf{Field} & \textbf{Type} & \textbf{Req} & \textbf{Default} & \textbf{Description} \\
\midrule
\texttt{id} & \texttt{id} & Yes & --- & Unique section identifier \\
\texttt{label} & \texttt{string} & Yes & --- & Section title \\
\texttt{time\_range} & \texttt{TimeRange?} & No & inherit & Override document time range \\
\texttt{tracks} & \texttt{list<ref<Track>>?} & No & all & Subset of tracks to show \\
\texttt{description} & \texttt{string?} & No & \texttt{null} & Section description \\
\texttt{page\_break} & \texttt{bool} & No & \texttt{false} & Hint for paginated output \\
\bottomrule
\end{tabular}
\caption{Section fields}
\label{tab:section-fields}
\end{table}

%% ============================================================================
%% LEGEND
%% ============================================================================
\subsection{Legend}
\label{sec:ir-legend}

The legend documents the meaning of statuses, categories, and visual conventions used
in the timeline. If omitted, the renderer may auto-generate a legend from used values.

\begin{lstlisting}[language=yaml,caption={Legend example}]
legend:
  show: true
  position: bottom-right
  entries:
    - key: in-progress
      label: In Progress
      description: Currently being worked on
    - key: at-risk
      label: At Risk
      description: May miss planned date
    - key: launch
      label: Launch Event
      category: true
\end{lstlisting}

\begin{table}[h]
\centering
\small
\begin{tabular}{llccp{5cm}}
\toprule
\textbf{Field} & \textbf{Type} & \textbf{Req} & \textbf{Default} & \textbf{Description} \\
\midrule
\texttt{show} & \texttt{bool} & No & \texttt{true} & Whether to render legend \\
\texttt{position} & \texttt{string?} & No & theme & Position hint \\
\texttt{title} & \texttt{string?} & No & \texttt{"Legend"} & Legend title \\
\texttt{entries} & \texttt{list<LegendEntry>} & No & auto & Legend entries \\
\bottomrule
\end{tabular}
\caption{Legend fields}
\label{tab:legend-fields}
\end{table}

\begin{table}[h]
\centering
\small
\begin{tabular}{llccp{5cm}}
\toprule
\textbf{Field} & \textbf{Type} & \textbf{Req} & \textbf{Default} & \textbf{Description} \\
\midrule
\texttt{key} & \texttt{string} & Yes & --- & Status or category key \\
\texttt{label} & \texttt{string} & Yes & --- & Human-readable label \\
\texttt{description} & \texttt{string?} & No & \texttt{null} & Extended description \\
\texttt{category} & \texttt{bool} & No & \texttt{false} & True if this is a category (not status) \\
\bottomrule
\end{tabular}
\caption{LegendEntry fields}
\label{tab:legend-entry-fields}
\end{table}


%% ============================================================================
%% ID AND REFERENCE SYSTEM
%% ============================================================================
\subsection{ID and Reference System}
\label{sec:ir-refs}

\subsubsection{Identifier Format}

All \texttt{id} fields must be:
\begin{itemize}
  \item Non-empty strings
  \item Composed of lowercase letters, digits, and hyphens
  \item Starting with a letter
  \item Unique within their entity type AND globally within the document
\end{itemize}

\textbf{Regex:} \verb|^[a-z][a-z0-9-]*$|

\textbf{Rationale for global uniqueness:} References may appear in contexts where the
entity type is ambiguous (e.g., annotation targets). Global uniqueness eliminates
the need for type prefixes in references.

\begin{lstlisting}[language=yaml,caption={ID examples}]
# Valid IDs
id: api-v2
id: q3-release
id: mobile-team-ios

# Invalid IDs
id: API-v2         # uppercase
id: 2026-release   # starts with digit
id: api_v2         # underscore
\end{lstlisting}

\subsubsection{References}

A reference is a string containing the \texttt{id} of another entity. The IR uses
bare strings for references (not structured objects) for YAML terseness:

\begin{lstlisting}[language=yaml]
activities:
  - id: feature-a
    track: platform      # ref<Track>
    group: engineering   # ref<Group>

annotations:
  - type: callout
    target: feature-a    # ref<Activity>
\end{lstlisting}

\textbf{Invariant:} Every reference must resolve to exactly one entity in the document.
A validator must reject documents with dangling references.

\subsubsection{Reference Namespacing}

Because IDs are globally unique, references are unambiguous without type prefixes.
The schema context determines what types are valid:

\begin{itemize}
  \item \texttt{activity.track}: must reference a Track
  \item \texttt{annotation.target}: may reference Activity or Milestone
  \item \texttt{group.children}: may reference Track or Group
\end{itemize}

A validator checks that the referenced entity has the correct type for its context.

%% ============================================================================
%% WELL-FORMEDNESS AND INVARIANTS
%% ============================================================================
\subsection{Well-Formedness Rules}
\label{sec:ir-invariants}

A Timeline IR document is \textbf{well-formed} if and only if all the following
invariants hold:

\subsubsection{Structural Invariants}

\begin{enumerate}
  \item \textbf{Version present:} \texttt{version} field exists and matches a known
        specification version.
  \item \textbf{Required fields:} All fields marked ``Req: Yes'' are present and non-null.
  \item \textbf{At least one track:} The \texttt{tracks} list is non-empty.
  \item \textbf{Unique IDs:} All \texttt{id} values are globally unique within the document.
  \item \textbf{Valid ID format:} All \texttt{id} values match \verb|^[a-z][a-z0-9-]*$|.
\end{enumerate}

\subsubsection{Referential Invariants}

\begin{enumerate}[resume]
  \item \textbf{References resolve:} Every reference points to an existing entity.
  \item \textbf{Reference types match:} References point to entities of the expected type
        (e.g., \texttt{activity.track} references a Track, not a Group).
  \item \textbf{No circular groups:} The group parent-child graph is acyclic.
\end{enumerate}

\subsubsection{Temporal Invariants}

\begin{enumerate}[resume]
  \item \textbf{Valid time range:} If \texttt{metadata.time\_range} has both start and end,
        then \texttt{end >= start}.
  \item \textbf{Activity duration non-negative:} For activities with concrete start and end,
        \texttt{end >= start}.
  \item \textbf{Date format valid:} All date values conform to the date model
        (\S\ref{sec:date-model}).
  \item \textbf{Activity date-source exclusive:} Every activity must satisfy exactly one of:
        (a)~\texttt{span} is present and neither \texttt{start} nor \texttt{end} is present; or
        (b)~\texttt{start} is present and \texttt{span} is absent (\texttt{end} optional).
        Co-presence of \texttt{span} with \texttt{start} or \texttt{end} is a hard
        well-formedness error: \texttt{SPAN\_START\_CONFLICT}.
  \item \textbf{Date anchor present when required:} If any date field in the document
        contains \texttt{now} or a relative offset (e.g., \texttt{+3m}, \texttt{-2w}),
        then at least one of \texttt{metadata.today} or \texttt{metadata.created} must be
        present. Absence of both is a hard error: \texttt{DATE\_ANCHOR\_MISSING}.
  \item \textbf{Track index values unique:} When \texttt{track.index} is present, all
        supplied values must be unique non-negative integers within the document.
\end{enumerate}

\subsubsection{Value Invariants}

\begin{enumerate}[resume]
  \item \textbf{Progress in range:} If \texttt{progress} is present, $0 \le \text{progress} \le 1$.
  \item \textbf{Status valid:} All \texttt{status} values are members of the Status enum.
  \item \textbf{Axis unit valid:} If present, \texttt{axis\_unit} is a member of AxisUnit enum.
\end{enumerate}

\subsubsection{Soft Warnings (Valid but Suspicious)}

The following are not errors but may indicate authorial mistakes:

\begin{itemize}
  \item Activity or milestone outside \texttt{metadata.time\_range} (will be clipped)
  \item \texttt{progress: 1.0} with \texttt{status: in-progress} (likely stale)
  \item Empty \texttt{activities} and \texttt{milestones} lists (vacuous timeline)
  \item Unused tracks (no activities reference them)
\end{itemize}


%% ============================================================================
%% EXAMPLES
%% ============================================================================
\subsection{Complete Examples}
\label{sec:ir-examples}

The following examples demonstrate realistic Timeline IR documents. Each is a complete,
valid document showcasing different timeline use cases.

\subsubsection{Example 1: Product Roadmap}

A quarterly product roadmap with multiple teams, showing how the IR represents a
typical software planning artifact.

\begin{lstlisting}[language=yaml,caption={Product roadmap example},label={lst:roadmap-example}]
version: "1.0"

metadata:
  title: "Product Roadmap 2026"
  subtitle: "Engineering & Design"
  author: "Product Team"
  created: 2026-06-09
  time_range:
    start: 2026-Q1
    end: 2026-Q4
  axis_unit: quarter
  theme: corporate-blue

tracks:
  - id: platform
    label: Platform
    description: Core infrastructure and APIs

  - id: mobile
    label: Mobile Apps
    description: iOS and Android applications

  - id: design
    label: Design System
    description: Component library and guidelines

groups:
  - id: foundation
    label: Foundation Work
    members: [api-v2, component-lib]

activities:
  - id: api-v2
    label: API v2 Migration
    track: platform
    start: 2026-Q1
    end: 2026-Q2
    status: in-progress
    progress: 0.6
    category: infrastructure

  - id: mobile-redesign
    label: Mobile App Redesign
    track: mobile
    start: 2026-Q2
    end: 2026-Q3
    status: planned
    description: Complete visual refresh

  - id: component-lib
    label: Component Library
    track: design
    start: 2026-Q1
    end: 2026-Q4
    status: in-progress
    progress: 0.3

  - id: analytics
    label: Analytics Integration
    track: platform
    span: 2026-Q3
    status: planned

milestones:
  - id: api-ga
    label: API v2 GA
    date: 2026-06-30
    track: platform
    status: planned
    icon: flag

  - id: app-launch
    label: App Store Launch
    date: 2026-Q4
    track: mobile
    status: planned
    category: launch

annotations:
  - type: today-marker
    date: now

legend:
  show: true
  entries:
    - key: infrastructure
      label: Infrastructure
      category: true
    - key: launch
      label: Launch Event
      category: true
\end{lstlisting}

\subsubsection{Example 2: Executive Program Timeline}

A high-level transformation program timeline suitable for board presentations,
demonstrating status communication and risk visibility.

\begin{lstlisting}[language=yaml,caption={Executive program timeline},label={lst:exec-timeline}]
version: "1.0"

metadata:
  title: "Digital Transformation Program"
  subtitle: "FY26 Executive View"
  author: "PMO"
  created: 2026-06-09
  time_range:
    start: 2026-01-01
    end: 2026-12-31
  axis_unit: month
  theme: executive

tracks:
  - id: strategy
    label: Strategy & Governance

  - id: technology
    label: Technology Modernization

  - id: people
    label: People & Change

activities:
  - id: strategy-define
    label: Strategy Definition
    track: strategy
    start: 2026-01-01
    end: 2026-03-31
    status: done
    progress: 1.0

  - id: vendor-selection
    label: Vendor Selection
    track: technology
    start: 2026-02-01
    end: 2026-04-30
    status: done

  - id: platform-build
    label: Platform Build
    track: technology
    start: 2026-04-01
    end: 2026-09-30
    status: in-progress
    progress: 0.35

  - id: data-migration
    label: Data Migration
    track: technology
    start: 2026-07-01
    end: 2026-11-30
    status: at-risk
    description: Dependency on legacy system documentation

  - id: training
    label: Training Program
    track: people
    start: 2026-06-01
    end: 2026-12-31
    status: planned

  - id: change-mgmt
    label: Change Management
    track: people
    start: 2026-03-01
    end: ongoing
    status: in-progress

milestones:
  - id: board-approval
    label: Board Approval
    date: 2026-03-15
    status: done
    icon: star

  - id: go-live
    label: Go Live
    date: 2026-10-01
    track: technology
    status: planned
    icon: flag

  - id: program-close
    label: Program Close
    date: 2026-12-15
    status: planned

annotations:
  - type: today-marker
    date: now

  - type: callout
    target: data-migration
    text: "Risk: Legacy documentation gaps"
    position: above

  - type: bracket
    targets: [platform-build, data-migration]
    label: "Critical Path"

sections:
  - id: h1
    label: "H1 2026: Foundation"
    time_range:
      start: 2026-01-01
      end: 2026-06-30

  - id: h2
    label: "H2 2026: Execution"
    time_range:
      start: 2026-07-01
      end: 2026-12-31
\end{lstlisting}


\subsubsection{Example 3: Architecture Evolution Timeline}

A technical architecture timeline showing system evolution over multiple years,
demonstrating coarse granularity and technical categorization.

\begin{lstlisting}[language=yaml,caption={Architecture evolution timeline},label={lst:arch-timeline}]
version: "1.0"

metadata:
  title: "Platform Architecture Evolution"
  subtitle: "2026-2028 Technical Roadmap"
  author: "Architecture Team"
  created: 2026-06-09
  time_range:
    start: 2026
    end: 2028
  axis_unit: half
  theme: technical

tracks:
  - id: frontend
    label: Frontend

  - id: backend
    label: Backend Services

  - id: data
    label: Data Platform

  - id: infra
    label: Infrastructure

activities:
  - id: react-migration
    label: React 19 Migration
    track: frontend
    start: 2026-H1
    end: 2026-H2
    status: in-progress
    category: modernization

  - id: micro-frontend
    label: Micro-Frontend Architecture
    track: frontend
    start: 2027-H1
    end: 2027-H2
    status: planned
    category: architecture

  - id: api-gateway
    label: API Gateway
    track: backend
    span: 2026-H1
    status: done
    category: infrastructure

  - id: service-mesh
    label: Service Mesh Adoption
    track: backend
    start: 2026-H2
    end: 2027-H1
    status: planned
    category: infrastructure

  - id: event-driven
    label: Event-Driven Architecture
    track: backend
    start: 2027
    end: 2028
    status: tentative
    category: architecture

  - id: lakehouse
    label: Data Lakehouse
    track: data
    start: 2026-H1
    end: 2027-H1
    status: in-progress
    progress: 0.4
    category: data

  - id: ml-platform
    label: ML Platform
    track: data
    start: 2027
    end: ongoing
    status: tentative
    category: ai

  - id: k8s-migration
    label: Kubernetes Migration
    track: infra
    span: 2026
    status: in-progress
    progress: 0.7
    category: infrastructure

  - id: multi-cloud
    label: Multi-Cloud Strategy
    track: infra
    start: 2027
    end: 2028
    status: planned
    category: infrastructure

milestones:
  - id: legacy-sunset
    label: Legacy System Sunset
    date: 2027-H1
    track: backend
    status: planned
    icon: flag

  - id: cloud-native
    label: Cloud Native Milestone
    date: 2026
    track: infra
    status: planned

annotations:
  - type: period
    start: 2026-H2
    end: 2027-H1
    label: "Transition Period"
    style: highlight

legend:
  entries:
    - key: modernization
      label: Modernization
      category: true
    - key: architecture
      label: Architecture Change
      category: true
    - key: infrastructure
      label: Infrastructure
      category: true
    - key: data
      label: Data Platform
      category: true
    - key: ai
      label: AI/ML
      category: true
\end{lstlisting}


%% ============================================================================
%% EXTENSIBILITY
%% ============================================================================
\subsection{Extensibility}
\label{sec:ir-extensibility}

The IR is designed to be extended without breaking existing consumers.

\subsubsection{Metadata Extension}

Every entity type includes a \texttt{metadata} field (type \texttt{map<string,any>})
for application-specific data:

\begin{lstlisting}[language=yaml]
activities:
  - id: feature-x
    label: Feature X
    track: platform
    span: 2026-Q2
    metadata:
      jira_key: PROJ-1234
      owner: alice@example.com
      priority: high
\end{lstlisting}

Renderers should ignore unrecognized metadata keys. Source integrations may use
metadata to preserve round-trip fidelity with external systems.

\subsubsection{Custom Categories}

The \texttt{category} field accepts arbitrary strings. The theme engine maps
categories to visual styles. New categories can be introduced without IR changes:

\begin{lstlisting}[language=yaml]
activities:
  - id: security-audit
    category: security   # custom category
    # ...

legend:
  entries:
    - key: security
      label: Security Work
      category: true
\end{lstlisting}

\subsubsection{Version Compatibility}

The \texttt{version} field enables forward compatibility:

\begin{itemize}
  \item Consumers should accept documents with higher minor versions (1.1, 1.2, etc.)
        and ignore unknown fields.
  \item Major version changes (2.0) may introduce breaking changes; consumers should
        reject documents with unsupported major versions.
\end{itemize}

%% ============================================================================
%% SERIALIZATION
%% ============================================================================
\subsection{Serialization and Syntax}
\label{sec:ir-serialization}

\subsubsection{Canonical Format}

YAML 1.2 is the canonical surface syntax for Timeline IR documents. Design decisions
prioritizing YAML:

\begin{itemize}
  \item \textbf{Minimal punctuation:} No braces for simple mappings, no quotes for
        most strings.
  \item \textbf{Line-oriented:} Each field on its own line for clean diffs.
  \item \textbf{Stable ordering:} Fields should appear in specification order for
        consistency; validators may enforce this.
  \item \textbf{Comments preserved:} YAML allows comments; tools should preserve them.
\end{itemize}

\subsubsection{Alternative Formats}

Tooling may support additional formats with lossless conversion:

\begin{itemize}
  \item \textbf{JSON:} Direct mapping; verbose but universally parseable.
  \item \textbf{TOML:} Alternative for users who prefer it.
\end{itemize}

The IR is defined abstractly; the YAML syntax is one serialization, not the definition.

\subsubsection{Git Friendliness}

For minimal diff churn:

\begin{enumerate}
  \item Lists use block style (one item per line), not flow style.
  \item Fields appear in consistent order.
  \item Auto-generated fields (like \texttt{updated} timestamps) are optional and
        may be omitted in version-controlled documents.
  \item IDs are stable slugs, not generated UUIDs.
\end{enumerate}

\begin{lstlisting}[language=yaml,caption={Git-friendly style}]
# Good: block style, stable order
activities:
  - id: feature-a
    label: Feature A
    track: platform
    span: 2026-Q2

# Avoid: flow style, hard to diff
activities: [{id: feature-a, label: Feature A, track: platform, span: 2026-Q2}]
\end{lstlisting}

%% ============================================================================
%% VALIDATION
%% ============================================================================
\subsection{Validation}
\label{sec:ir-validation}

A conforming validator must check all invariants in \S\ref{sec:ir-invariants}.
Validation is a function:

\[
\text{validate} : \text{Document} \to \text{Valid} \mid \text{Errors}
\]

\subsubsection{Error Reporting}

Validation errors should include:
\begin{itemize}
  \item Path to the offending field (e.g., \texttt{activities[2].track})
  \item The invariant violated
  \item The actual value
  \item Suggested fix (when determinable)
\end{itemize}

\subsubsection{Strict vs. Lenient Mode}

\begin{itemize}
  \item \textbf{Strict mode:} All invariants enforced; required for rendering.
  \item \textbf{Lenient mode:} Warnings instead of errors for soft issues; useful
        for work-in-progress documents.
\end{itemize}

%% ============================================================================
%% RELATIONSHIP TO OTHER COMPONENTS
%% ============================================================================
\subsection{Relationship to Other Components}
\label{sec:ir-relationships}

The IR is the stable contract between pipeline stages:

\begin{description}
  \item[Ingestion (upstream)] Source adapters (ADO, GitHub, Markdown, etc.) produce
        IR documents. The IR defines the target shape; adapters handle mapping.

  \item[Theme Engine (downstream)] Themes consume IR and produce styled visual
        specifications. Themes interpret \texttt{status}, \texttt{category}, and
        hint fields (\texttt{color}, \texttt{icon}) but cannot require additional
        IR fields.

  \item[Renderer (downstream)] The renderer consumes themed IR (or IR + theme) and
        produces output (SVG, PNG, PDF, PPTX). The IR remains rendering-agnostic;
        pixel-level decisions are not IR concerns.

  \item[Agent Integration] AI agents generate IR directly. The IR's explicit typing,
        clear required/optional fields, and forgiving structure enable reliable
        generation. Agents need not understand rendering.
\end{description}

\textbf{Contract stability:} Downstream components may not require IR fields beyond
this specification. Extension fields (\texttt{metadata}) are pass-through.


%% ============================================================================
%% BUILD VS ADOPT: SURVEY OF EXISTING REPRESENTATIONS
%% (Research: David, 2026-06-10)
%% ============================================================================
\subsection{Build vs.\ Adopt: A Survey of Existing Representations}
\label{sec:ir-build-vs-adopt}

Before committing to a fully greenfield IR, this section surveys existing formats,
grammars, and standards as potential adoption or extension candidates. Every
candidate is assessed against the five criteria derived from the design principles
(\S\ref{sec:ir-structure}):
\begin{enumerate}
  \item \textbf{Visual-communication focus:} roadmap/milestone/swimlane semantics,
        \emph{not} task-scheduling (no dependencies, resource levelling, or
        critical path);
  \item \textbf{Git-friendly:} line-oriented text, stable field order, minimal diff
        churn;
  \item \textbf{LLM-generatable:} clear required/optional schema, simple references,
        forgiving structure;
  \item \textbf{Render/theme separation:} clean boundary between \emph{what} is
        communicated and \emph{how} it looks; and
  \item \textbf{Open licence with community or specification body.}
\end{enumerate}

%% ------------------------------------------------------------------
\subsubsection{Timeline Text Languages}
%% ------------------------------------------------------------------

\paragraph{Markwhen.}
Markwhen~\cite{markwhen} (MIT licence) is the closest existing tool to our
surface-syntax goals. Events are written as \texttt{date: label} or
\texttt{start/end: label} lines; \texttt{group:} blocks provide swimlane-like
structure; \texttt{\#tag} tokens allow lightweight categorisation. The format
is text-native, line-oriented, and git-diffable.

Critical gaps exist, however. There is no first-class status concept
(no equivalent of \texttt{status: at-risk}); granularity is limited to ISO
dates and informal text strings (no \texttt{2026-Q3} or fiscal-quarter
shorthand); render/theme separation is absent---the view container interprets
colour from tags, not a separate theme layer; static SVG, PDF, and PPTX output
are unsupported; and Markwhen does not publish a versioned, machine-verifiable
schema for its internal JSON parse tree, which makes reliable LLM generation
fragile. The editor (markwhen.com) is proprietary; only the parser and view
container are open source.

\paragraph{Mermaid timeline and Gantt.}
The Mermaid \texttt{timeline} type~\cite{mermaiddocs2024} groups events into
sections on a free-text time axis. It has no swimlane concept, no status enum,
no PPTX output, and no render/theme separation (styles are hard-coded in the
renderer). The Mermaid \texttt{gantt} type~\cite{mermaidgantt2024} adds
dependency keywords (\texttt{crit}, \texttt{done}, \texttt{active})---precisely
the scheduling semantics this specification rejects. Neither type exposes a
stable IR that downstream tooling can consume.

\paragraph{PlantUML Gantt.}
PlantUML's \texttt{@startgantt}~\cite{plantumlgantt} is experimental and
scheduling-oriented (dependency declarations, work-remaining percentages). Its
verbose, UML-rooted syntax is not LLM-friendly for presentation timelines.

\paragraph{D2 and Structurizr.}
D2~\cite{d2lang} is a general-purpose diagram DSL with no native timeline type.
Structurizr~\cite{structurizr} is domain-specific to software architecture (C4
model). Neither is a relevant adoption candidate.

%% ------------------------------------------------------------------
\subsubsection{Visualization Grammars}
%% ------------------------------------------------------------------

\paragraph{Vega-Lite.}
Vega-Lite~\cite{vegalite2017} is the most significant analogy in this
category. It demonstrates that a high-level declarative JSON grammar can
cleanly separate \emph{what to visualize} (data, encodings, marks) from
\emph{how to render it} (scales, axes, guides, legends), with actual rendering
delegated to the Vega runtime. This is precisely the WHAT/HOW philosophy that
this IR embodies (see \S\ref{sec:ir-relationships}).

Vega-Lite is, however, a statistical-graphics grammar, not a roadmap grammar.
Encoding a swimlane timeline in Vega-Lite requires manual mark composition:
a \texttt{rect} mark with \texttt{y: track}, \texttt{x: start},
\texttt{x2: end}, plus separate layers for milestones and annotations. There is
no first-class concept of \texttt{tracks}, \texttt{milestones}, \texttt{status},
\texttt{span}, or \texttt{annotations}. The resulting spec is verbose
(100{+} JSON lines for a five-activity roadmap), fragile for LLM generation,
and not human-editable in the YAML sense intended here. Vega-Lite is the
strongest \emph{architectural analogy}---its declarative WHAT/HOW separation is
the direct design precedent for this IR---but cannot be adopted wholesale.

\paragraph{The Grammar of Graphics and ggplot2.}
Wilkinson's grammar~\cite{grammarofgraphics2005} and Wickham's layered
refinement~\cite{layeredgrammar2010} establish the theoretical foundation for
declarative visual grammars: decompose a graphic into data, aesthetics,
geometry, statistics, and scales. The key insight---separating semantic content
from visual treatment---directly motivates this IR's separation of
tracks/activities/status from theme/colours/layout. These are
\emph{conceptual donors}, not adoptable formats.

\paragraph{Observable Plot and Apache ECharts.}
Observable Plot (ISC licence~\cite{observable-plot}) is a concise JavaScript
charting API. It has no native swimlane roadmap type and requires imperative
programming, not declarative text authoring. Apache ECharts has a
\texttt{timeline} component, but it is an animated-playback control that cycles
through multiple chart states over time---not a swimlane roadmap grammar.
Neither is suitable for adoption.

%% ------------------------------------------------------------------
\subsubsection{Timeline Data Models}
%% ------------------------------------------------------------------

\paragraph{Knight Lab TimelineJS.}
TimelineJS~\cite{timelinejs} accepts a flat JSON array of events with
\texttt{start\_date}, \texttt{end\_date}, \texttt{text.headline}, and
\texttt{text.text}. It is storytelling-oriented (narrative slides, embedded
media) with no swimlanes, no status enum, no PPTX output, and no render/theme
separation.

\paragraph{vis-timeline.}
vis-timeline~\cite{vistimeline} is a browser-only interactive widget requiring
JavaScript item objects (\texttt{\{id, content, start, end, group\}}). There is
no declarative text format, no static SVG/PDF/PPTX export, and no theme-layer
concept. Not adoptable.

\paragraph{react-chrono and Timeline Storyteller.}
react-chrono~\cite{react-chrono} (MIT) accepts a JavaScript items array with
\texttt{title}, \texttt{cardTitle}, \texttt{cardDetailedText}; date fields are
plain strings with no semantic granularity and no swimlanes. Microsoft
Research's Timeline Storyteller~\cite{timeline-storyteller} (open-sourced 2019)
accepts a simple \texttt{[\{start, end, label\}]} JSON array---no swimlanes,
no status, no maintained community. Neither is adoptable.

%% ------------------------------------------------------------------
\subsubsection{Temporal and Calendar Standards}
%% ------------------------------------------------------------------

\paragraph{iCalendar (RFC 5545).}
RFC~5545~\cite{ical-rfc5545} defines \textsc{vevent} with properties
\texttt{DTSTART}, \texttt{DTEND}, \texttt{SUMMARY}, \texttt{DESCRIPTION},
\texttt{CATEGORIES}, \texttt{STATUS} (\texttt{TENTATIVE}, \texttt{CONFIRMED},
\texttt{CANCELLED}), and \texttt{URL}. This vocabulary is the outstanding
\emph{donor}: Mark's \texttt{start}/\texttt{end}/\texttt{label}/\texttt{description}/
\texttt{category}/\texttt{status}/\texttt{url} all map directly to iCalendar
terms, establishing prior-art legitimacy for these names
(see Table~\ref{tab:alignment}).

The \textsc{ics} wire format is wholly unsuitable: property-folded lines
(\texttt{DTSTART:20260609T140000Z}), no YAML/JSON structure, and no swimlane,
visual-status, or render-pipeline concept. The \texttt{STATUS} enum covers
scheduling state (\texttt{CONFIRMED}) but not visual-communication urgency
(\texttt{at-risk}, \texttt{blocked}). \textbf{Not adoptable; strongest single
vocabulary donor.}

\paragraph{Schema.org/Event.}
The schema.org \texttt{Event} type~\cite{schemaorg-event} defines
\texttt{name}, \texttt{startDate}, \texttt{endDate}, \texttt{description},
\texttt{url}, and \texttt{eventStatus}. These are the canonical web vocabulary
for events and corroborate Mark's field naming. Schema.org is a
semantic-markup standard (JSON-LD/Microdata), not a visual representation
format. \textbf{Not adoptable; useful vocabulary donor.}

\paragraph{W3C OWL-Time.}
OWL-Time~\cite{owltime} (W3C Recommendation, 2022) formalises temporal
entities as \texttt{Instant} and \texttt{Interval} with object properties
\texttt{hasBeginning}, \texttt{hasEnd}, and \texttt{hasTemporalDuration}.
Crucially, omitting \texttt{hasEnd} explicitly represents an open/ongoing
interval---the precise semantics of the IR's \texttt{end: ongoing}. OWL-Time
implements Allen's thirteen interval relations~\cite{allen1983} and is the
authoritative formal reference for the IR's open-ended activity semantics.

\paragraph{Allen's Interval Algebra.}
Allen's 1983 paper~\cite{allen1983} establishes 13 binary relations between
time intervals: \emph{before}, \emph{meets}, \emph{overlaps}, \emph{starts},
\emph{during}, \emph{finishes}, \emph{equals}, and their converses. The
algebra provides formal vocabulary for the uncertain/open date model: an
activity with \texttt{end: ongoing} participates in the \emph{started-by}
relation with the document's \texttt{time\_range} without requiring a concrete
end date. Not a format; the formal grounding for temporal semantics.

\paragraph{ISO 8601.}
ISO~8601~\cite{iso8601} is already the foundation of the date model
(\S\ref{sec:date-model}). The IR's \texttt{2026-06-09}, \texttt{2026-06}, and
\texttt{2026} forms are ISO~8601; the \texttt{2026-Q2} and \texttt{2026-H1}
extensions follow the same \texttt{YYYY-unit} pattern as a natural extension.
No further alignment is needed.

%% ------------------------------------------------------------------
\subsubsection{Scheduling and PM Formats}
%% ------------------------------------------------------------------

TaskJuggler~\cite{taskjuggler} (GPL~v2), MS Project XML~\cite{msprojectxml},
and Primavera P6 are scheduling grammars built around \emph{computed} dates:
dependency resolution, resource levelling, and critical-path analysis. All are
explicitly out of scope per the Timeline Grammar decision. MS Project XML
additionally fails the git-friendly criterion: the format is non-deterministic
between saves (GUIDs and timestamps change on every write). No vocabulary from
these formats merits borrowing beyond what iCalendar and schema.org already
provide.

%% ------------------------------------------------------------------
\subsubsection{Analogous Document IRs}
%% ------------------------------------------------------------------

The Pandoc AST~\cite{pandoc} is the strongest structural analogy: a typed,
versioned, language-agnostic JSON document IR with a \texttt{pandoc-api-version}
key, a \texttt{meta} block, and a sequence of typed block elements. This
pattern---\texttt{version} + \texttt{metadata} + typed entity lists---maps
directly onto the root structure of \S\ref{sec:ir-structure}. The Dragon
Book~\cite{dragonbook2006} provides theoretical backing for a typed,
pipeline-stable IR as the contract between compilation stages. Both are
\emph{structural analogies}, not adoptable formats.

%% ------------------------------------------------------------------
\subsubsection{Comparison Table}
%% ------------------------------------------------------------------

Table~\ref{tab:build-vs-adopt} scores the seven strongest candidates against
the five criteria. \textbf{Y}~=~good fit; \textbf{P}~=~partial;
\textbf{N}~=~poor fit or not applicable.

\begin{table}[h]
\centering
\small
\begin{tabular}{p{3.2cm}ccccp{1.6cm}c}
\toprule
\textbf{Candidate}
  & \textbf{Vis.}
  & \textbf{Git}
  & \textbf{LLM}
  & \textbf{Sep.}
  & \textbf{Licence}
  & \textbf{Tooling} \\
\midrule
Markwhen~\cite{markwhen}
  & Y & Y & P & N & MIT & P \\
Mermaid timeline~\cite{mermaiddocs2024}
  & P & Y & P & N & MIT & Y \\
iCalendar VEVENT~\cite{ical-rfc5545}
  & N & N & N & N & IETF & Y \\
TimelineJS JSON~\cite{timelinejs}
  & Y & P & Y & N & GPL & P \\
vis-timeline~\cite{vistimeline}
  & Y & N & N & N & MIT & Y \\
Vega-Lite~\cite{vegalite2017}
  & P & P & P & Y & BSD-3 & Y \\
TaskJuggler~\cite{taskjuggler}
  & N & Y & N & N & GPLv2 & P \\
\bottomrule
\end{tabular}
\caption{Build-vs-Adopt scoring. Vis.~=~visual-communication focus (not
scheduling); Git~=~git-friendly text format; LLM~=~LLM-generatable schema;
Sep.~=~render/theme separation; Tooling~=~active community. Y~=~good;
P~=~partial; N~=~poor.}
\label{tab:build-vs-adopt}
\end{table}

No candidate scores \textbf{Y} on all five dimensions. Markwhen---the
closest---fails on render/theme separation, lacks a stable machine-verifiable
schema, and has no PPTX output path.

%% ------------------------------------------------------------------
\subsubsection{Recommendation: \textsc{Build} with Borrowed Vocabulary}
%% ------------------------------------------------------------------

\textbf{Verdict: \textsc{Build} a bespoke IR. No existing format is
adoptable wholesale and no viable extension profile exists.}

Two orthogonal gaps eliminate every candidate:

\begin{enumerate}
  \item \textbf{Semantic gap.} All formats with real communities (Mermaid,
        iCalendar, Vega-Lite) are scheduling-oriented, statistical-graphics-%
        oriented, or calendar-oriented. None is a visual-communication roadmap
        grammar. Defining a ``profile'' of iCalendar or a ``mark recipe'' in
        Vega-Lite would require adding swimlanes, visual status, coarse-grain
        dates, PPTX output, and a theme engine as first-class concepts---at
        which point we have built a new layer regardless.

  \item \textbf{Pipeline gap.} The requirement of a static, multi-backend
        render pipeline (SVG / PDF / PPTX) with a separate theme engine has no
        counterpart in any existing open-source tool. Mermaid, Markwhen, and
        vis-timeline all bake rendering into the format. Vega-Lite delegates to
        a JavaScript runtime, making it unsuitable for PPTX or print-PDF output
        without adaptation that would negate any adoption benefit.
\end{enumerate}

A \textsc{build} decision is therefore justified. It must not be a
wheel-reinvention: the following \textbf{vocabulary and semantic donors} are
explicitly leveraged so that the IR speaks standard language wherever standards
exist.

\begin{description}
  \item[ISO 8601~\cite{iso8601}]
        Date string syntax (\texttt{2026-06-09}, \texttt{2026-Q2}).
        Already used in the date model (\S\ref{sec:date-model});
        no change needed.

  \item[iCalendar RFC~5545~\cite{ical-rfc5545}]
        Field naming: \texttt{DTSTART}~$\to$~\texttt{start},
        \texttt{DTEND}~$\to$~\texttt{end},
        \texttt{SUMMARY}~$\to$~\texttt{label},
        \texttt{DESCRIPTION}~$\to$~\texttt{description},
        \texttt{CATEGORIES}~$\to$~\texttt{tags}/\texttt{category},
        \texttt{URL}~$\to$~\texttt{url}.
        Status vocabulary: \texttt{TENTATIVE} appears verbatim in the Status
        enum; \texttt{CANCELLED} maps to \texttt{cancelled}.

  \item[Schema.org/Event~\cite{schemaorg-event}]
        Confirms \texttt{name}/\texttt{startDate}/\texttt{endDate}/
        \texttt{description}/\texttt{url} as the canonical web vocabulary for
        events. Mark's \texttt{label}/\texttt{start}/\texttt{end} are direct
        equivalents (shorter names preferred for human authoring).

  \item[W3C OWL-Time~\cite{owltime} and Allen's Algebra~\cite{allen1983}]
        Open-interval semantics: omitting \texttt{hasEnd} denotes an ongoing
        interval, validating \texttt{end: ongoing}. The Instant/Interval
        duality formalises the \texttt{Milestone} (instant) vs.\
        \texttt{Activity} (interval) distinction.

  \item[Vega-Lite~\cite{vegalite2017}]
        Architectural precedent: strict WHAT/HOW separation between a
        declarative specification (this IR) and the rendering runtime
        (Theme Engine + Renderer). The pattern of a versioned, typed YAML/JSON
        document as the sole contract between source and output is adopted from
        Vega-Lite's design.

  \item[Markwhen~\cite{markwhen} and Mermaid~\cite{mermaiddocs2024}]
        Surface-syntax precedents: readable event lines, section/group
        structure, tag-based categorisation. These inform future higher-level
        authoring syntaxes that compile down to this IR.

  \item[Pandoc AST~\cite{pandoc}]
        Structural analogy: versioned, typed, language-agnostic document IR with
        \texttt{version} + \texttt{meta} + typed entity lists. Validates the
        root structure (\S\ref{sec:ir-structure}) and the principle of a
        stable, machine-verifiable schema.
\end{description}

%% ------------------------------------------------------------------
\subsubsection{Alignment of Mark's IR with Prior Art}
%% ------------------------------------------------------------------

Table~\ref{tab:alignment} maps Mark's field names to their closest standard
equivalents. Fields marked \emph{Original} have no adequate standard analog
and are justified by the visual-communication use case.

\begin{table}[h]
\centering
\small
\begin{tabular}{p{2.4cm}p{3.0cm}p{2.0cm}p{4.0cm}}
\toprule
\textbf{IR field} & \textbf{iCal / schema.org} & \textbf{OWL-Time} & \textbf{Assessment} \\
\midrule
\texttt{start}, \texttt{end}
  & \texttt{DTSTART/DTEND}; \texttt{startDate/endDate}
  & \texttt{hasBeginning/} \texttt{hasEnd}
  & \textbf{Aligned.} Standard names; shorter preferred for authoring. \\
\texttt{label}
  & \texttt{SUMMARY}; \texttt{name}
  & ---
  & \textbf{Aligned.} Shorter form preferred for YAML. \\
\texttt{description}
  & \texttt{DESCRIPTION}
  & ---
  & \textbf{Aligned.} \\
\texttt{url}
  & \texttt{URL}; \texttt{url}
  & ---
  & \textbf{Aligned.} \\
\texttt{tags}
  & \texttt{CATEGORIES} (multi)
  & ---
  & \textbf{Aligned.} Correct plural form. \\
\texttt{category}
  & \texttt{CATEGORIES} (single)
  & ---
  & \textbf{Aligned.} \\
\texttt{status: tentative}
  & \texttt{STATUS:TENTATIVE}
  & ---
  & \textbf{Aligned.} Verbatim iCalendar value. \\
\texttt{status: cancelled}
  & \texttt{STATUS:CANCELLED}
  & ---
  & \textbf{Aligned.} Lower-case per IR convention. \\
\texttt{status: at-risk}
  & ---
  & ---
  & \textbf{Original.} No standard analog; justified by executive-presentation idiom. \\
\texttt{status: blocked}
  & ---
  & ---
  & \textbf{Original.} No standard analog; justified by visual-communication need. \\
\texttt{end: ongoing}
  & ---
  & omit \texttt{hasEnd}
  & \textbf{Semantically aligned.} Explicit encoding is clearer than omission. \\
\texttt{span}
  & ---
  & ---
  & \textbf{Original.} Convenient shorthand; no standard equivalent. \\
\texttt{track}
  & \texttt{CALENDAR} (weak)
  & ---
  & \textbf{Original.} Markwhen uses \texttt{group}; \texttt{track} is clearer. \\
\texttt{progress}
  & ---
  & ---
  & \textbf{Original.} No standard analog; no change needed. \\
\texttt{Activity} (type)
  & \textsc{vevent} (duration $>$ 0)
  & \texttt{Interval}
  & \textbf{Aligned.} Correct modelling of time-spanning entities. \\
\texttt{Milestone} (type)
  & \textsc{vevent} (zero dur.)
  & \texttt{Instant}
  & \textbf{Aligned.} Correct modelling of point-in-time markers. \\
\bottomrule
\end{tabular}
\caption{Alignment of Mark's IR field names and types with prior-art standards.}
\label{tab:alignment}
\end{table}

\paragraph{Flags for Mark and Leslie (no schema changes required).}
\begin{itemize}
  \item \textbf{No schema changes are recommended.} Mark's field choices are
        well-aligned with iCalendar RFC~5545, schema.org/Event, and OWL-Time.
        The original contributions (\texttt{at-risk}, \texttt{blocked},
        \texttt{span}, \texttt{progress}, \texttt{track}) are justified by the
        visual-communication use case and have no adequate standard equivalent.

  \item \textbf{Document vocabulary donors in the spec preamble.} This section
        serves that purpose. Future contributors should not propose renames (e.g.,
        \texttt{name} for \texttt{label}, \texttt{startDate} for \texttt{start})
        without consulting this survey, since shorter names are deliberate and
        standards-corroborated.

  \item \textbf{Optional surface alias \texttt{type: event}:} schema.org uses
        \texttt{Event} and OWL-Time uses \texttt{Instant} for point-in-time
        items. Exposing \texttt{type: event} as a surface-syntax alias for
        milestones could ease LLM generation without changing the IR schema.
        Flagged for future surface-language design.
\end{itemize}

